\section{什么是 Good Eval：Outcome Validity 与失效模式}

有了任务集合，下一问题是分数是否真的代表目标能力。课程采用 outcome validity 作为核心：评估结果应与任务意图、真实成功和可接受过程一致。一个稳定、可重复但测错目标的 benchmark，仍然是坏评估。

\subsection{Outcome validity：测到“真正完成”，而不是替代信号}

slides 40--41 用 Agent 执行链说明，用户请求经 Agent、工具和环境产生结果，evaluation system 必须判断该结果是否满足意图。若只评分最终回复，可能忽略环境没有改变；若只检查状态，可能忽略回复误导用户；若只看 reward，又可能遗漏成本和违规。

\lecturefigure{slides-images/slide-040.png}{官方 slide 40：Good eval 的 Agent--environment--evaluator 结构}
\lecturefigure{slides-images/slide-041.png}{官方 slide 41：Outcome validity 在完整执行链中的位置}

\begin{importantbox}{Outcome validity 的三个一致性}
意图一致：评分规则对应用户真正目标；状态一致：环境中的事实变化与声明一致；过程一致：达到目标没有违反权限、安全或资源约束。三者至少需要分别观测。Agent “说已完成”不是 outcome，工具返回成功码也不一定代表业务状态正确。
\end{importantbox}

\subsection{文本、代码、状态变更与多步推理的不同判据}

slides 42--45 分别给出四类 outcome。文本结果可检查内容、依据和指令满足；代码生成需构建、测试、静态分析与人工质量；环境状态变更要做 state matching 和约束检查；多步推理既看最终答案，也要防止 shortcut、随机猜中和过程违规。不同 outcome 需要不同 oracle，不能统一成 LLM judge。

\lecturefigure{slides-images/slide-042.png}{官方 slide 42：文本结果的内容、依据与约束判断}
\lecturefigure{slides-images/slide-043.png}{官方 slide 43：代码生成的构建、测试与质量验证}
\lecturefigure{slides-images/slide-044.png}{官方 slide 44：环境状态变更的 state matching}
\lecturefigure{slides-images/slide-045.png}{官方 slide 45：多步推理的答案、过程与一致性验证}

\begin{knowledgebox}{读图：Outcome-specific oracle}
文本任务可使用 claim verification、引用检查和 rubric；代码任务使用编译、unit/integration tests、lint、安全扫描；状态任务比较数据库/DOM/API 后态与目标不变量；推理任务可用最终答案、proof checker、步骤约束和重复采样。应把强确定性 oracle 放在主评分，将主观 judge 限于可读性或解释等维度。
\end{knowledgebox}

\begin{warningbox}{测试通过并不证明实现正确}
测试覆盖有限、环境可能被修改、随机性可能碰巧成功。代码 Agent benchmark 应隔离测试、保护 verifier、检查 diff 范围并运行 hidden tests；状态任务要验证无关字段没有被破坏；多步任务要检查重复执行是否幂等。Outcome validity 是证据组合，不是单一绿灯。
\end{warningbox}

\subsection{Ways Eval Can Go Wrong}

slide 46 汇总数据噪声/偏差、任务不真实、数据不平衡、shortcut、评估成本与不可靠 judge 等失效模式。更隐蔽的问题是 Goodhart's law：当分数成为优化目标，系统会寻找规则漏洞。Agent 有工具和长轨迹，投机空间比静态模型更大。

\lecturefigure{slides-images/slide-046.png}{官方 slide 46：数据、任务、shortcut、成本与 judge 失效模式}

\begin{knowledgebox}{评估失效审计表}
\begin{longtable}{p{2.8cm}p{4.2cm}p{4.2cm}}
\toprule
\textbf{风险} & \textbf{证据} & \textbf{缓解} \\
\midrule
数据偏差 & 分桶表现、来源分布 & 重采样、补充真实长尾 \\
污染 & 重合检索、异常熟悉度 & 私有/动态题、时间切分 \\
Shortcut & 高分轨迹异常短或违规 & hidden checks、过程审计 \\
Judge 偏差 & 人类校准、位置/长度敏感 & 随机化、多 judge、oracle \\
成本失控 & token、工具、wall time 分布 & 预算上限、成本前沿 \\
\bottomrule
\end{longtable}
\end{knowledgebox}

\subsection{本章小结}

Good eval 的中心不是指标数量，而是 outcome validity。文本、代码、状态和推理需要不同证据；数据、judge 和 harness 都可能失效。下一章用 benchmark case study 检查这些原则如何落地。

